% Cloudra 0.1.0rc1 软件使用说明书（软著申请用），tectonic 直接编译
\PassOptionsToPackage{fontset=none}{ctex}
\documentclass[zihao=-4,a4paper,fontset=none]{ctexart}
\usepackage[top=25mm,bottom=25mm,left=25mm,right=25mm]{geometry}
\usepackage{ctex}
\usepackage{booktabs}
\usepackage{array}
\usepackage{amsmath}
\emergencystretch=2em
\usepackage{enumitem}
\usepackage{fancyhdr}
\usepackage{hyperref}
\hypersetup{hidelinks}
\usepackage{caption}
\captionsetup{labelsep=quad}
\usepackage{xcolor}
\usepackage{listings}
\setmonofont{Menlo}
\setCJKmonofont{Songti SC}
\lstset{
  basicstyle=\ttfamily\small,
  frame=single,
  rulecolor=\color{gray!60},
  breaklines=true,
  keepspaces=true,
  aboveskip=8pt,
  belowskip=8pt,
}
\pagestyle{fancy}
\fancyhf{}
\fancyhead[L]{\small Cloudra 0.1.0rc1 软件使用说明书}
\fancyhead[R]{\small 第 \thepage{} 页}
\renewcommand{\headrulewidth}{0.4pt}

\title{\bfseries Cloudra 软件使用说明书\\[6pt] \large 0.1.0rc1}
\author{Cloudra Team}
\date{2026年9月23日}

\setCJKmainfont{Songti SC}
\begin{document}

\maketitle
\thispagestyle{empty}

\tableofcontents

\section{软件概述}

\subsection{软件名称与版本}
\begin{itemize}[leftmargin=4em,itemsep=2pt]
  \item 软件全称：Cloudra自主有机小分子机理计算Agent软件
  \item 软件简称：Cloudra
  \item 版本号：0.1.0rc1
  \item 开发完成日期：2026年9月23日
\end{itemize}

\subsection{软件用途}
Cloudra 是一款面向均相有机小分子反应机理研究的自主计算 Agent 软件。系统承袭「职责清晰、先算后核」的设计，以事件驱动运行时协调总控 Agent 与领域子 Agent：从反应物／产物等化学信息出发提出机理假设，自动完成结构准备、过渡态初猜与搜索、虚频确认、反应路径连通及热化学等环节；以机理图承载物种状态与转化关系，并配套几何／电子结构感知、计算任务调度、危险操作审批与过程材料留存，支撑可复查的机理解析流程。
\subsection{运行环境}
\begin{itemize}[leftmargin=4em,itemsep=2pt]
  \item 操作系统：Linux、macOS
  \item 编程语言：Python 3.11 及以上
  \item 依赖环境：OpenAI 兼容的大语言模型 API
  \item 界面环境：现代浏览器（Chrome / Edge / Firefox）
\end{itemize}

\section{软件功能模块}

\subsection{事件驱动运行时内核}
软件的核心是事件驱动的运行时内核（runtime\_v2），由以下组件构成：
\begin{itemize}[itemsep=2pt]
  \item \textbf{事件日志与事件总线}：全部会话行为（消息、工具调用、审批、提醒等）以事件形式顺序记录并广播，前端通过 SSE 实时订阅，完整日志可持久化重放。
  \item \textbf{命令总线}：所有写操作先事件后状态，保证多 Agent 并发下的状态一致性；Agent 的生命周期由状态机管理，消息提交与唤醒按统一投递策略执行。
  \item \textbf{Agent 会话状态机}：每个 Agent 拥有独立身份与状态（idle、running、waiting\_approval、waiting\_input、finished 等），会话上下文按事件持久化，重启后恢复已提交的投影；运行中的任务按恢复规则收口。
  \item \textbf{并发运行器}：每个 Agent 一个异步任务，真正并行执行且可取消；后台任务统一登记并受取消与关闭流程管理；工具调用生命周期持久化。
  \item \textbf{固定轮次管道}：每轮依次经过权限、无人值守、方法、DAG、目标、计划、知识与引导等策略阶段；LLM 输出以 token 增量流式推送，每轮完整推理过程落持久化事件。
\end{itemize}

\subsection{多 Agent 协作体系}
\begin{itemize}[itemsep=2pt]
  \item \textbf{Coordinator（协调者）}：负责接收用户需求、拆解任务并调度各专业子 Agent。
  \item \textbf{八类专业子 Agent}：见表~\ref{tab:agents}。
  \item \textbf{子 Agent 调度与蜂群}：支持按模板批量展开的 Swarm 并发调度（可设并发上限）、同步等待与父任务取消级联。
  \item \textbf{AgentHub 论坛}：Agent 间公告板，支持发帖、浏览、检索与跨 Agent 来电呼叫。
  \item \textbf{统一信箱}：用户与 Agent 间消息统一投递，支持定向回复与呼叫。
  \item \textbf{提醒调度}：拦截定时工具调用，由调度机制在到期时投递提醒并唤醒对应 Agent。
\end{itemize}

\begin{table}[htbp]
  \centering
  \caption{专业化 Agent 矩阵}
  \label{tab:agents}
  \begin{tabular}{ll>{\raggedright\arraybackslash}p{8.5cm}}
    \toprule
    Agent 名称 & 标识 & 功能定位 \\
    \midrule
    Coordinator & a-0 & 任务协调与总控 \\
    Explorer & a-1 & 项目文件结构导航与发现 \\
    Writer & a-2 & 文档撰写、报告总结、代码编写 \\
    Researcher & a-3 & 学术文献检索与专业数据调研 \\
    small\_mol & a-4 & 小分子建模、构象与结构文件处理 \\
    md & a-5 & 分子动力学模拟、力场与轨迹分析 \\
    dft & a-6 & 量子化学计算（ORCA／Gaussian／xTB 等）输入输出与流程编排 \\
    experiment & a-7 & 实验数据处理与操作辅助 \\
    env & a-8 & 系统环境配置与软件安装管理 \\
    \bottomrule
  \end{tabular}
\end{table}

\subsection{推理与上下文管理}
\begin{itemize}[itemsep=2pt]
  \item \textbf{流式输出}：LLM token 增量以事件流直达前端，逐字渲染；瞬时数据不落日志、不参与重放。
  \item \textbf{持久化推理记录}：每轮完整思考过程落一条持久事件，历史轮次刷新后仍可回看。
  \item \textbf{常驻与一次性注入}：目标、DAG 等常驻说明以固定注入方式参与每轮组装；通知类一次性注入在轮内统一收集。
  \item \textbf{手动上下文压缩}：对超长对话可手动或按预算触发结构化摘要压缩，在轮边界安全执行。
  \item \textbf{长期记忆}：MEMORY.md 长期记忆文件自动加载，跨会话保留研究偏好与结论。
\end{itemize}

\subsection{工具套件}
\begin{itemize}[itemsep=2pt]
  \item \textbf{Chemtool 化学工具}：结构构建与 SMILES 转换、量子化学输入编写与日志解析（能量、频率、虚频判别等）、光谱与格式转换；集成过渡态初猜（插值／路径搜索）、结构与电子结构感知，减少对手写三维坐标的依赖。
  \item \textbf{自动过渡态搜索管线}：将结构准备、原子映射、路径初猜、校验、鞍点优化、频率与 IRC 判定等步骤封装为可自动推进的流程；Agent 启动后跟踪进度，并在异常断点处调整后再续跑。
  \item \textbf{Searchtool 检索工具}：arXiv、OpenAlex 等学术文献检索，以及原生联网搜索适配。
  \item \textbf{Bashtool 命令工具}：POSIX Shell 执行，内置命令分级、文件变更审批与危险命令守卫。
  \item \textbf{Writetool 写作工具}：结构化文档、实验记录与报告写入。
  \item \textbf{Plantool 计划工具}：研究计划与机理图操作，提供机理草稿、证据绑定等图操作指令。
  \item \textbf{Managetool 管理工具}：子 Agent 调度、蜂群展开、AgentHub 论坛、定时提醒。
  \item \textbf{Skilltool 技能工具}：加载科研方法技能包（第一性原理计算、文献综述、论文写作等）。
  \item \textbf{Stoptool 终止工具}：按目标完成度门禁终止会话，防止过早收尾。
\end{itemize}

\subsection{机理图引导编排与自动过渡态搜索}
软件支持按会话开启的机理图（DAG）引导模式，用于均相有机反应机理解析：
\begin{itemize}[itemsep=2pt]
  \item \textbf{化学信息机理图}：以物种节点与转化关系承载化学状态与约束，模型提出机理草稿后经接受落为交付图。
  \item \textbf{自动过渡态搜索}：关键基元步骤可走声明驱动的搜索管线，自动完成初猜、优化与路径确认；失败时给出可操作反馈并由 Agent 调整后续跑。
  \item \textbf{计算复核}：声明为证据源的工具调用完成后进入复核流程，对虚频数目、路径连通等结果做检查后再绑定到对应步骤。
  \item \textbf{交付检查}：结束任务时检查机理图覆盖与关键证据是否齐备，未完成则提示继续，避免过早交卷。
  \item \textbf{督导模式}：可选科学督导角色，用于接受／驳回草稿并对关键绑定做复核。
\end{itemize}

\subsection{安全与权限系统}
\begin{itemize}[itemsep=2pt]
  \item \textbf{五级权限模型}：READ（读）、WRITE（写）、MANAGE（管理子 Agent）、SCRIPT（执行脚本）、DESTRUCTIVE（破坏性操作）。
  \item \textbf{危险操作审批流}：危险工具调用挂起并等待用户确认；请求与决策全部事件化、可追溯。
  \item \textbf{YOLO 自动审批守卫}：可选全自动模式，由独立 LLM 守卫对审批请求进行安全判断。
  \item \textbf{专项守卫}：坐标守卫（拦截绕过审查直接改写量子化学坐标文件的写法）、命令习惯守卫（提示 heredoc/sed -i 等易错写文件命令）、委派守卫（提示 Coordinator 自行提交计算任务）。
\end{itemize}

\subsection{计算执行与数据管理}
\begin{itemize}[itemsep=2pt]
  \item \textbf{本地与 Slurm 计算}：独立调度守护从中央计算库认领就绪任务，本地拉起监视进程持续解析输出；支持多节点注册表与 Slurm 作业提交、轮询与续租。
  \item \textbf{SQLite 持久化}：会话级数据库承载事件日志、DAG 图存储与复核收件箱，支持事件重放与会话上下文投影恢复。
  \item \textbf{不可变产物快照}：工具产生的文件以内容寻址快照保存，支持短期下载授权。
  \item \textbf{备份恢复}：内置 backup 子命令，在线快照 + SHA-256 清单，支持校验与恢复。
\end{itemize}

\section{软件安装与配置}

\subsection{安装步骤}
\begin{lstlisting}[language=bash]
# 克隆代码仓库
git clone <仓库地址>
cd <仓库目录>

# 安装（依赖统一声明在 pyproject.toml；可选安装化学扩展）
pip install -e .
pip install -e ".[chem]"   # 可选：RDKit、nmrglue 等化学依赖
\end{lstlisting}

\subsection{初始化配置}
\begin{lstlisting}[language=bash]
cloudra config init    # 交互式初始化：LLM 地址、模型、密钥等
cloudra config show    # 查看当前配置（敏感信息自动脱敏）
cloudra config set coordinator.model gpt-4o   # 设置单个配置项
cloudra config get coordinator.model           # 读取单个配置项
\end{lstlisting}

\section{软件使用方法}

\subsection{启动图形界面}
\begin{lstlisting}[language=bash]
cloudra                 # 构建并启动 Web 服务（默认打开浏览器）
cloudra web --host 127.0.0.1   # 指定本机绑定地址
\end{lstlisting}

前端界面由三部分组成：
\begin{itemize}[itemsep=2pt]
  \item \textbf{项目栏与 Agent 栏}：切换会话项目、查看各 Agent 状态（空闲/运行/等待审批/等待输入）。
  \item \textbf{对话区}：流式对话输出、推理过程分段折叠、工具调用行、Markdown 与 KaTeX 公式渲染。
  \item \textbf{检查器（Inspector）}：管理目标、机理图编排、权限、审批请求与上下文。
\end{itemize}

\subsection{常用管理指令}
\begin{lstlisting}[language=bash]
cloudra --version      # 查看版本
cloudra stats tokens   # 查看 LLM token 用量
cloudra scheduler      # 前台运行计算调度守护
cloudra stop           # 停止服务与守护进程
\end{lstlisting}

\section{典型使用场景}

\subsection{反应机理研究（机理图模式）}
用户在项目设置中开启机理图引导，向 Coordinator 提出反应机理解析目标（例如环加成或轴手性相关扭转能垒）。系统引导提出机理草稿并落图；对关键步骤启动自动过渡态搜索管线，计算完成后进行虚频与路径连通复核，再将结果绑定到对应节点／边，最终交付完整机理图与能垒信息。

\subsection{量子化学计算流程}
用户描述目标分子与计算需求，Coordinator 调度 dft／small\_mol 等子 Agent：生成合法输入并提交计算，由监视进程解析输出，提取能量、轨道能级、振动频率等指标并写入报告；涉及过渡态时优先走自动搜索管线而非仅手写输入文件。

\subsection{文献调研与综述}
用户要求调研某课题最新进展，Researcher 子 Agent 检索 arXiv 与 OpenAlex 数据库，Writer 子 Agent 依据检索结果撰写综述，产出结构与引用规范的 Markdown 报告。

\section{技术支持}
\begin{itemize}[leftmargin=4em,itemsep=2pt]
  \item 开发团队：Cloudra Team
  \item 反馈渠道：项目仓库 Issue 页面
\end{itemize}

\vfill
\begin{center}
\small 本说明书适用于 Cloudra 0.1.0rc1 版本软件
\end{center}

\end{document}
